\pdfinfoomitdate=1
\pdftrailerid{}
\pdfsuppressptexinfo=15
\newcommand{\DoNotLoadEpstopdf}{}
\documentclass[11pt]{article}
\usepackage[T1]{fontenc}
\usepackage[utf8]{inputenc}
\usepackage{lmodern,amsmath,amssymb,amsthm,mathtools,booktabs,array}
\usepackage[margin=1in]{geometry}
\usepackage{microtype}
\usepackage[hidelinks]{hyperref}
\usepackage{enumitem}
\setlist{itemsep=3pt,topsep=5pt}
\allowdisplaybreaks[2]
\emergencystretch=2em
\newtheorem{theorem}{Theorem}[section]
\newtheorem{lemma}[theorem]{Lemma}
\newtheorem{proposition}[theorem]{Proposition}
\newtheorem{corollary}[theorem]{Corollary}
\theoremstyle{definition}\newtheorem{definition}[theorem]{Definition}
\theoremstyle{remark}\newtheorem{remark}[theorem]{Remark}
\newcommand{\OO}{\mathcal O}
\newcommand{\NN}{\mathbb N}
\newcommand{\EE}{\mathbb E}
\newcommand{\Var}{\operatorname{Var}}
\newcommand{\odd}{\operatorname{odd}}
\newcommand{\stir}[2]{\genfrac{[}{]}{0pt}{}{#1}{#2}}
\newcommand{\ii}{\mathrm i}
\newcommand{\dd}{\,\mathrm d}
\DeclareMathOperator{\Rea}{Re}
\DeclareMathOperator{\Ima}{Im}
\hypersetup{pdftitle={Report185 Decorated logarithmic eta products and OEIS A330499},pdfauthor={Research report},pdfsubject={Uniform all orders oscillatory coefficient asymptotics and threshold inverses}}
\title{Report185\\[5pt]Decorated logarithmic eta products\\[3pt]and OEIS A330499}
\author{Uniform oscillatory asymptotics and threshold inverses}
\date{3 October 2026}
\begin{document}
\maketitle
\begin{abstract}
We prove an all-fixed-orders coefficient theorem for the composition
$F_B(z)=\sum_{k\ge1}\log(1+B(z)^k)$, where $B$ is analytic with nonnegative
coefficients and its occupied sizes have difference span one. If
$B(\rho)=1$ strictly inside the disk of analyticity, the leading coefficient
is $\pi^2/(12\mu)$ after multiplication by $\rho^n$, where $\mu$ is the
mean of the decorating size law. The first correction is
$n^{-3/4}H_0(\sqrt n)$, followed by half-power orders; each $H_j$ is an
explicit entire trigonometric series, bounded and uniformly almost periodic
on the real axis. Its mode damping is determined by the variance divided
by the squared mean. A pole-subtracted boundary integral and a uniform
stationary-phase estimate justify summing the actual remainders, even
though accumulating root-of-unity cusps prevent standard Delta-domain
transfer. For $B(z)=-\log(1-z)$, we identify the exact leading constant
of A330499, give its first three oscillatory terms, prove signed
subsequences that exclude a pure $1/n$ refinement, and derive
rounding-aware inverse envelopes to every fixed order. Exact finite
checks and optional floating-point diagnostics accompany the proof.
\end{abstract}

\paragraph{Scope.}
The reusable contribution is the decorated-product coefficient theorem,
including its uniform remainder. The eta transformation, stationary
phase, labelled supernecklaces and inverse reversion are classical
inputs. The bounded source comparison in Section~\ref{sec:sources} is
not a worldwide priority determination. All orders and the decoration
are fixed before $n\to\infty$. The constants in the asymptotic remainder
are existence constants; this report does not supply a certified
large-$n$ cutoff or a certified numerical inverse algorithm.

\newpage
{\small\tableofcontents}
\newpage

\section{Main results and their meaning}\label{sec:main}
Let
\begin{equation}\label{eq:model}
 L(z)=-\log(1-z),\qquad
 F(z)=\sum_{k\ge1}\log(1+L(z)^k)
     =\sum_{n\ge0}a_n\frac{z^n}{n!}.
\end{equation}
All logarithms in the defining series are the power-series logarithms
at the origin. The sequence $(a_n)_{n\ge0}$ is OEIS A330499~\cite{oeis}.
Its first values are
\[
 0,\ 1,\ 2,\ 13,\ 71,\ 558,\ 5344,\ 60926,\ 766898,\ 10759096.
\]
The OEIS record inspected on 3 October 2026 gives the leading equivalent
$n!c(1-e^{-1})^{-n}$ with the displayed decimal $c=0.478656\ldots$.
Define
\begin{equation}\label{eq:logconstants}
 \rho=1-e^{-1},\qquad M=e-1,\qquad C=\frac{\pi^2}{12M}
       =0.478656655620749637061291495659\ldots.
\end{equation}
The decimal is a numerical evaluation of the exact expression, not an
interval certificate.

\begin{theorem}[A330499 at every fixed order]\label{thm:A}
There are real entire functions $H_j$, $j\ge0$, bounded on the real axis,
such that for every fixed integer $N\ge1$,
\begin{equation}\label{eq:Aall}
 a_n=n!\rho^{-n}\left\{C+
       \sum_{j=0}^{N-1}n^{-3/4-j/2}H_j(\sqrt n)
       +\OO_N(n^{-3/4-N/2})\right\}.
\end{equation}
Each $H_j$ and each of its derivatives is given by an absolutely
convergent trigonometric series. The first three functions are explicit
in Section~\ref{sec:three}; a finite recipe for every further function
is given in Theorem~\ref{thm:general}. In particular,
\begin{equation}\label{eq:threeintro}
 \frac{a_n\rho^n}{n!}=C+n^{-3/4}H_0(\sqrt n)
        +n^{-5/4}H_1(\sqrt n)+n^{-7/4}H_2(\sqrt n)
        +\OO(n^{-9/4}).
\end{equation}
\end{theorem}

The oscillation in \eqref{eq:Aall} is a power-scale effect, not an
exponentially small correction. In particular, there is no constant
$c_1$ for which $a_n\rho^n/n!=C+c_1/n+\OO(n^{-2})$.
Section~\ref{sec:sign} proves this from an exact rational first-harmonic
tail estimate and explicit signed subsequences. A leading equivalent
by itself does not reveal this obstruction.

The proof is most useful in a broader setting. For a nonnegative
analytic decoration $B$, the numbers $b_j\rho^j$ at $B(\rho)=1$ define
a probability law. If its mean and variance are $\mu$ and $v$, the
$m$th oscillation has damping
\begin{equation}\label{eq:damping}
 \exp\left(-\frac{\pi^2mv}{\mu^2}\right).
\end{equation}
This identifies a structural smoothing parameter. A small variance
relative to the squared mean makes many modes relevant; it does not
justify taking a zero-variance limit of a theorem whose constants depend
on the decoration.

\paragraph{Why the proof uses the coefficient boundary.}
The modular eta identity is exact, but the coefficient boundary
approaches its singular point tangentially. The local real-axis
expansion loses the oscillatory information. Moreover, singular
preimages of odd-order roots of unity accumulate at $\rho$ from
nearly vertical directions. Section~\ref{sec:cusps} proves that these
are actual nonremovable cusps. The standard Delta-domain hypothesis
for singularity analysis is therefore unavailable. We instead subtract
the dominant pole, use dominated boundary coefficient extraction, and
control every dual eta mode uniformly before summing its error.

\section{Exact coefficients and labelled interpretation}\label{sec:exact}
For an integer $m\ge1$, let $\odd(m)$ be its largest odd divisor,
and let $\sigma(m)=\sum_{d\mid m}d$. Put
\begin{equation}\label{eq:cm}
 c_m=\frac{\sigma(\odd(m))}{m}.
\end{equation}
The distinction between $c_m$ and the leading constant $C$ will be used
throughout.

\begin{lemma}[The positive logarithmic product coefficients]\label{lem:cm}
For $|q|<1$,
\begin{equation}\label{eq:F0}
 F_0(q):=\sum_{k\ge1}\log(1+q^k)=\sum_{m\ge1}c_mq^m.
\end{equation}
In particular, $0<c_m\le1+\log m$.
\end{lemma}
\begin{proof}
The defining logarithmic series converges locally normally. Writing
$P(q)=\prod_{k\ge1}(1-q^k)$ and using its analytic logarithm gives
\[
 F_0(q)=\log P(q^2)-\log P(q),\qquad
 \log P(q)=-\sum_{m\ge1}\sigma_{-1}(m)q^m,
\]
where $\sigma_{-1}(m)=\sum_{d\mid m}d^{-1}=\sigma(m)/m$.
Consequently the coefficient is
\[
 \sigma_{-1}(m)-\mathbf1_{2\mid m}\sigma_{-1}(m/2).
\]
Writing $m=2^\nu u$ with $u$ odd shows that this difference equals
$\sigma(u)/m$. Positivity follows, while
$c_m\le\sigma_{-1}(m)\le\sum_{j=1}^m1/j\le1+\log m$.
Alternatively, expanding the individual logarithms gives the exact
identity $c_m=\sum_{j\mid m}(-1)^{j+1}/j$.
\end{proof}

Unsigned Stirling numbers of the first kind are denoted by
$\stir nk$. They satisfy
\[
 \frac{L(z)^k}{k!}=\sum_{n\ge k}\stir nk\frac{z^n}{n!},
 \qquad
 \stir{n+1}k=n\stir nk+\stir n{k-1}.
\]
Thus \eqref{eq:F0} yields the finite exact transform
\begin{equation}\label{eq:stirling}
 a_n=\sum_{k=1}^n\stir nk\,(k-1)!\,\sigma(\odd(k)),\qquad a_0=0.
\end{equation}
This is an integer formula with positive weights. It is also an
independent way to evaluate the coefficients without summing truncated
analytic products numerically.

\begin{proposition}[Marked cycles of cycles and strict growth]\label{prop:counts}
The number $a_n$ counts labelled cycles of cycles with the following
additional choice: if there are $k$ inner cycles, choose an odd divisor
$d$ of $k$ and one of $d$ colors. Furthermore,
\begin{equation}\label{eq:monotone}
 a_{n+1}>n a_n\quad(n\ge1),
\end{equation}
and $(a_n)_{n\ge0}$ is strictly increasing.
\end{proposition}
\begin{proof}
There are $\stir nk$ ways to decompose $n$ distinct labels into $k$
nonempty cycles. These $k$ distinguishable cycles can be arranged into
an outer directed cycle in $(k-1)!$ ways. The additional finite choice
has $\sum_{d\mid k,\ d\text{ odd}}d=\sigma(\odd(k))$ possibilities.
This gives \eqref{eq:stirling}. The underlying unmarked labelled
cycle-of-cycles construction is classical supernecklace
combinatorics~\cite{supernecklace,ac}; the divisor-colored version here
is the direct interpretation of the displayed weight.

Write $w_k=(k-1)!\sigma(\odd(k))>0$. Applying the Stirling recurrence
and taking $\stir nk=0$ outside $0\le k\le n$ gives
\[
 a_{n+1}=n a_n+\sum_{j=0}^n\stir nj w_{j+1}>n a_n.
\]
The last sum is positive for $n\ge1$. Together with $a_0=0$ and
$a_1=1$, this proves the assertion.
\end{proof}

Unmarked supernecklaces have generating function $-\log(1-L(z))$
and a different leading scale $n!\rho^{-n}/n$. The finite marking in
Proposition~\ref{prop:counts} is not asymptotically innocuous: the
logarithmic eta product has a pole-sized leading term and oscillatory
power corrections.

\section{The general analytic decoration theorem}\label{sec:general}
Assume that
\begin{equation}\label{eq:Bassump}
 B(z)=\sum_{j\ge1}b_jz^j,\qquad b_j\ge0,
 \qquad B\text{ holomorphic for }|z|<R,
 \qquad 0<\rho<R,\quad B(\rho)=1.
\end{equation}
The support is required to have \emph{difference span one}:
\begin{equation}\label{eq:span}
 \gcd\{j-k:b_jb_k>0\}=1.
\end{equation}
Here the gcd of the differences is the positive generator of the
subgroup they span. In particular the support is not a singleton.
This hypothesis is stronger than $\gcd\{j:b_j>0\}=1$.
For example, $(z+z^3)/2$ satisfies the latter condition, but
$B(-1)=-1$ gives another boundary contact.

Define the positive-integer-valued variable $K$ by
\begin{equation}\label{eq:law}
 \Pr\{K=j\}=b_j\rho^j,
 \qquad \mu=\EE K,\quad v=\Var K>0.
\end{equation}
Since $\rho<R$, this law has a moment generating function analytic
near zero. Denote its cumulants by $\kappa_r$, so
$\kappa_1=\mu$, $\kappa_2=v$. Let
\begin{equation}\label{eq:modes}
 F_B(z)=F_0(B(z)),\quad \alpha_m=2\pi^2m,
 \quad d=\frac{v}{2\mu^2},\quad
 q_1=\frac{\kappa_3}{6\mu^2}-\frac{v^2}{4\mu^3}.
\end{equation}
The dual-mode parameter $\alpha_m$ is unrelated to the sequence value $a_n$.

On a neighborhood of real $\theta=0$, use the logarithm branch with
$-\log B(\rho)=0$, and define the removable germ
\begin{equation}\label{eq:W}
 W(\theta)=\frac1{-\log B(\rho e^{\ii\theta})}
                  -\frac{\ii}{\mu\theta},
 \qquad E_\alpha(\theta)=e^{-\alpha(W(\theta)-d)}.
\end{equation}
The first terms are $W(\theta)=d+\ii q_1\theta+\OO(\theta^2)$.
Let
\begin{equation}\label{eq:Y}
 Y(w)=\left(\frac{w+\sqrt{w^2+4}}2\right)^2,
 \qquad Y(0)=1.
\end{equation}
The square root is the analytic branch positive at zero. The identity
$Y+Y^{-1}=2+w^2$ is exact.

\begin{theorem}[Uniform transfer at every fixed order]\label{thm:general}
Under \eqref{eq:Bassump}--\eqref{eq:span}, for every fixed integer
$N\ge1$,
\begin{equation}\label{eq:general}
 \rho^n[z^n]F_B(z)=\frac{\pi^2}{12\mu}
  +\sum_{j=0}^{N-1}n^{-3/4-j/2}H_j(\sqrt n)
  +\OO_{B,N}(n^{-3/4-N/2}).
\end{equation}
For $\alpha>0$, form the finite expression
\begin{equation}\label{eq:Qrecipe}
\begin{split}
 Q_j(\alpha)={}&\sum_{r+k=j}
 \frac{(-\ii)^r}{4^rr!k!}
 \left(\frac\alpha\mu\right)^{k/2}
 \left(\frac\mu\alpha\right)^{r/2}E_\alpha^{(k)}(0)\\[-2pt]
 &\hspace{25mm}\cdot
 \left.\frac{\mathrm d^{2r}}{\mathrm dw^{2r}}
                \{Y(w)^kY'(w)\}\right|_{w=0}.
\end{split}
\end{equation}
There is a real function $P_j$ with $Q_j=(-\ii)^jP_j$;
$P_j(\alpha)$ is $\alpha^{-j/2}$ times a polynomial in $\alpha$
of degree at most $2j$. Then
\begin{equation}\label{eq:Hrecipe}
\begin{split}
 H_j(u)=-\sum_{m\ge1}&c_m e^{-\alpha_md}
          \frac{(\alpha_m/\mu)^{1/4}}{\sqrt\pi}P_j(\alpha_m)\\
 &\cdot\cos\left(2u\sqrt{\alpha_m/\mu}+\frac\pi4+\frac{j\pi}2\right).
\end{split}
\end{equation}
Every such series, and each fixed derivative in $u$, converges
absolutely on the real axis and locally uniformly on the complex
plane. The functions $H_j$ are real entire and uniformly almost
periodic on the real axis; they and all their derivatives are bounded
there.
\end{theorem}

The adjective ``entire'' does not mean bounded on the complex plane.
Uniform almost periodicity means uniform approximation on the real
axis by finite trigonometric sums; it follows immediately from the
absolute convergence asserted in the theorem.

The first terms of the finite recipe are particularly simple:
\begin{equation}\label{eq:P01}
 P_0(\alpha)=1,\qquad
 P_1(\alpha)=\frac{q_1\alpha^{3/2}}{\sqrt\mu}
                      +\frac{3\sqrt\mu}{16\sqrt\alpha}.
\end{equation}
Thus
\begin{align}
 H_0(u)&=-\sum_{m\ge1}c_me^{-\alpha_md}
      \frac{(\alpha_m/\mu)^{1/4}}{\sqrt\pi}
      \cos\left(2u\sqrt{\alpha_m/\mu}+\frac\pi4\right),\label{eq:H0general}\\
 H_1(u)&=\sum_{m\ge1}c_me^{-\alpha_md}
      \frac{(\alpha_m/\mu)^{1/4}}{\sqrt\pi}P_1(\alpha_m)
      \sin\left(2u\sqrt{\alpha_m/\mu}+\frac\pi4\right).\label{eq:H1general}
\end{align}
In view of $\alpha_md=\pi^2mv/\mu^2$, these formulas give
\eqref{eq:damping} without fitting an empirical decay constant.

\subsection{Examples of admissible decorations}
The theorem applies to ordinary analytic functions as written; their
interpretation as labelled classes is optional. Table~\ref{tab:B}
shows common decorations. Each has strictly positive consecutive
coefficients, or support $\{1,2\}$, hence difference span one.
For the set decoration write $\ell=\log2$.
\begin{table}[htbp]
\centering
\small
\begin{tabular}{lllll}
\toprule
Decoration $B(z)$ & $\rho$ & $\mu$ & $v$ & $v/\mu^2$\\
\midrule
Lists $z/(1-z)$ & $1/2$ & $2$ & $2$ & $1/2$\\
Sets $e^z-1$ & $\ell$ & $2\ell$ & $2\ell(1-\ell)$ & $(1-\ell)/(2\ell)$\\
Cycles $-\log(1-z)$ & $1-e^{-1}$ & $M$ & $M$ & $1/M$\\
Two sizes $(z+z^2)/2$ & $1$ & $3/2$ & $1/4$ & $1/9$\\
\bottomrule
\end{tabular}
\caption{The variance ratio determines exponential mode damping.
Every parameter is exact.}\label{tab:B}
\end{table}
These entries follow from $\mu=\rho B'(\rho)$ and
$v=\rho^2B''(\rho)+\rho B'(\rho)-\mu^2$. A monomial decoration
has zero variance and is excluded. The theorem is not uniform in a
family approaching a monomial, changing its support span, or moving
$\rho$ to the boundary of analyticity.

\section{The exact eta identity and boundary extraction}\label{sec:boundary}
The proof of Theorem~\ref{thm:general} occupies this section and
Section~\ref{sec:uniform}. No unproved exchange of a fixed-mode
asymptotic expansion with an infinite sum is used.

\subsection{The transformed product}
The classical modular transformation of the Dedekind eta function
implies, for positive real $t$,
\begin{equation}\label{eq:Peta}
 P(e^{-t})=\sqrt{\frac{2\pi}{t}}
    \exp\left(\frac t{24}-\frac{\pi^2}{6t}\right)
    P(e^{-4\pi^2/t}).
\end{equation}
The logarithm and square-root branches extend analytically from
$t>0$ to $\Rea t>0$. See DLMF 23.18~\cite{eta} and the explicit
product identity (15) of Ardehali--Rosengren~\cite{ardehali}.
Subtracting the logarithm of \eqref{eq:Peta} from the same identity
with $t$ replaced by $2t$ gives
\begin{equation}\label{eq:etaexact}
 F_0(e^{-t})=\frac{\pi^2}{12t}-\frac{\log2}{2}+\frac t{24}
                 -\sum_{m\ge1}c_m e^{-2\pi^2m/t},\qquad \Rea t>0.
\end{equation}
Indeed, with $Q=e^{-2\pi^2/t}$, the remaining product logarithm is
$\log P(Q)-\log P(Q^2)=-F_0(Q)$, and $|Q|<1$ because
$\Rea(1/t)>0$. The identity holds first on the positive real axis and
then on the entire right half-plane by holomorphic continuation. This
also fixes the logarithmic branch without a possible hidden multiple
of $2\pi\ii$.

On the coefficient circle $z=\rho e^{\ii\theta}$, the cumulant identity
$\log B(\rho e^{\ii\theta})=\log\EE e^{\ii\theta K}$ gives
\begin{align}
 t(\theta):=-\log B(\rho e^{\ii\theta})
  &=-\ii\mu\theta+\frac v2\theta^2
           +\frac{\ii\kappa_3}{6}\theta^3+\OO(\theta^4),\label{eq:texpand}\\
 \frac1{t(\theta)}
  &=\frac{\ii}{\mu\theta}+d+\ii q_1\theta+\OO(\theta^2).
 \label{eq:invtexpand}
\end{align}
In particular $W$ in \eqref{eq:W} is analytic. Choose a sufficiently
small $\delta>0$ so that it is holomorphic on a neighborhood of
$[-\delta,\delta]$ and
\begin{equation}\label{eq:Wpositive}
 \Rea W(\theta)\ge d/2\quad (|\theta|\le\delta,\ \theta\in\mathbb R).
\end{equation}
This positive real part provides the summable mode damping on the
actual boundary.

\subsection{Why span one removes the other boundary contacts}
By the triangle inequality, $|B(\rho e^{\ii\theta})|\le1$.
Equality is possible exactly when all the unit complex numbers
$e^{\ii j\theta}$ for occupied sizes $j$ are identical. Equivalently,
$e^{\ii(j-k)\theta}=1$ for every pair of occupied sizes. The span-one
hypothesis is therefore exactly what implies
\begin{equation}\label{eq:strictcircle}
 |B(\rho e^{\ii\theta})|<1\quad
         (\theta\not\equiv0\pmod{2\pi}).
\end{equation}
On any compact arc avoiding zero, the upper bound is strictly below
one. Since $B$ is analytic beyond the circle, the defining
$F_0(B(z))$ series extends holomorphically through that arc. Hence its
boundary values are smooth there, with all derivatives bounded on
compact arcs.

\subsection{Subtracting the pole before taking boundary values}
Set $x=z/\rho$. A direct Taylor expansion at $x=1$ gives
\begin{equation}\label{eq:poleexpand}
 \frac1{-\log B(\rho x)}
 =\frac1{\mu(1-x)}+\frac{v-\mu}{2\mu^2}+\OO(1-x).
\end{equation}
For $|x|<1$,
$\Rea(1/(1-x))\ge1/2$. The constant contribution to the lower
bound in \eqref{eq:poleexpand} is
$1/(2\mu)+(v-\mu)/(2\mu^2)=d$. Thus, after shrinking the local
neighborhood,
\[
 \Rea\frac1{-\log B(\rho x)}\ge d/2
       \qquad (|x|<1,\ x\text{ near }1).
\]
The dual-mode sum in \eqref{eq:etaexact} is consequently bounded there
by $\sum_m c_m e^{-\alpha_md/2}<\infty$.

It follows that
\begin{equation}\label{eq:G}
 G(x)=F_B(\rho x)-\frac{\pi^2}{12\mu(1-x)}
\end{equation}
is bounded on the unit disk. Near $x=1$, this follows from
\eqref{eq:poleexpand}, the bounded mode sum and the analytic term
$t/24$; away from $1$, it follows from \eqref{eq:strictcircle} and
compactness. Boundary values exist and are smooth away from $1$,
and are bounded near it. Dominated convergence on radial circles
therefore gives Fourier coefficient extraction from those boundary
values. No assumption of analytic continuation across $1$ is made.

Choose a real even $C^\infty$ cutoff $\chi$, supported in
$(-\delta,\delta)$ and equal to one on a smaller neighborhood of
zero. The analytic pole-subtracted part of $\pi^2/(12t)$ and the
remaining elementary terms in \eqref{eq:etaexact}, multiplied by
$\chi$, are smooth functions of $\theta$. All terms outside the
cutoff are smooth as well. Repeated integration by parts makes their
$n$th Fourier coefficients $\OO_L(n^{-L})$ for every fixed $L$.
We have reduced the coefficient problem to
\begin{equation}\label{eq:reduction}
 \rho^n[z^n]F_B(z)-\frac{\pi^2}{12\mu}
 =-\sum_{m\ge1}c_mJ_{\alpha_m}(n)+\OO_L(n^{-L}),
\end{equation}
where
\begin{equation}\label{eq:J}
 J_\alpha(n)=\frac1{2\pi}\int_{-\delta}^{\delta}
   \chi(\theta)e^{-\ii(n\theta+\alpha/(\mu\theta))}
                     e^{-\alpha W(\theta)}\dd\theta.
\end{equation}
The value at the single point $\theta=0$ is irrelevant. The exchange
of the sum with the integral is justified by the summable majorant
$c_me^{-\alpha_md/2}$ from \eqref{eq:Wpositive}.

\section{Uniform stationary phase and the summed remainder}\label{sec:uniform}
The following estimate is the quantitative step needed for the
infinite mode sum. Its uniformity is stronger than a fixed-$m$
stationary-phase expansion.

\begin{lemma}[Uniform single-mode estimate]\label{lem:uniform}
Fix $\alpha_0>0$ and an integer $N\ge1$. There are positive constants
$\eta,c,C_N$ and a nonnegative integer $M_N$, depending on the fixed
decoration, $N$ and $\alpha_0$, such that for all sufficiently large $n$ and
$\alpha_0\le\alpha\le\eta n$,
\begin{equation}\label{eq:Jasym}
\begin{split}
 J_\alpha(n)={}&n^{-3/4}e^{-\alpha d}
      \frac{(\alpha/\mu)^{1/4}}{\sqrt\pi}
      \Rea\left\{e^{-\ii(2\sqrt{\alpha n/\mu}+\pi/4)}
                  \sum_{j=0}^{N-1}n^{-j/2}Q_j(\alpha)\right\}\\
 &+R_{N,\alpha}(n),
\end{split}
\end{equation}
with $Q_j$ as in \eqref{eq:Qrecipe} and
\begin{equation}\label{eq:uniformerror}
 |R_{N,\alpha}(n)|\le
 C_Nn^{-3/4-N/2}(1+\alpha)^{M_N}e^{-c\alpha}.
\end{equation}
The constants do not depend on $n$ or $\alpha$ in this range.
\end{lemma}

\begin{proof}
Put
\begin{equation}\label{eq:scaling}
 s=\sqrt{\frac\alpha{\mu n}},\qquad
 \lambda=\sqrt{\frac{\alpha n}{\mu}},\qquad \theta=sy.
\end{equation}
Then
\[
 n\theta+\frac\alpha{\mu\theta}
       =\lambda\phi(y),\qquad \phi(y)=y+y^{-1}.
\]
The phase has exactly two real stationary points, $y=1$ and $y=-1$,
and $\phi''(1)=2$, $\phi''(-1)=-2$. Choose $\eta$ small enough
that $s\le s_0$ and fixed neighborhoods of these two points lie in
$\chi(sy)=1$ for all allowed $s$. Use a fixed partition into those
two neighborhoods and the nonstationary complement.

\paragraph{Uniform amplitude bounds.}
On either fixed stationary neighborhood, every fixed-order derivative
in $y$ of $e^{-\alpha W(sy)}$ is bounded by
\begin{equation}\label{eq:ampbounds}
 C_j(1+\alpha)^{D_j}e^{-c\alpha},
\end{equation}
uniformly for $0<s\le s_0$. This follows from the chain rule, bounded
derivatives of $W$, and \eqref{eq:Wpositive}. Factors of $s$ are
bounded by fixed constants. The same estimate holds after an analytic
coordinate change with bounded derivatives on the chosen compact
neighborhood. Constants may change from line to line.

\paragraph{An explicit stationary-phase remainder bound.}
Near $y=1$, set $w=(y-1)/\sqrt y$, so $y=Y(w)$ and
$\phi(y)=2+w^2$. For a compactly supported smooth amplitude $f$, the
Gaussian integral gives
\begin{equation}\label{eq:gaussian}
 \int_{\mathbb R}e^{-\ii\lambda w^2}f(w)\dd w
 =e^{-\ii\pi/4}\sqrt{\frac\pi\lambda}
    \sum_{r=0}^{N-1}\frac{(-\ii)^r}{4^rr!\lambda^r}f^{(2r)}(0)
   +\mathcal R_N.
\end{equation}
For completeness, Fourier inversion expresses the left side as the
Gaussian factor times the integral of
$\widehat f(\xi)e^{\ii\xi^2/(4\lambda)}$. Taylor expansion of the
last exponential through degree $N-1$, using its real-axis remainder
bound, gives
\begin{equation}\label{eq:gaussianbound}
 |\mathcal R_N|\le C_N\lambda^{-N-1/2}
             \int_{\mathbb R}|\xi|^{2N}|\widehat f(\xi)|\dd\xi.
\end{equation}
The weighted Fourier norm is bounded by a fixed combination of the
$L^1$ norms of $f$ and its derivatives through order $2N+2$: use the
$L^1$ bound for $|\xi|\le1$ and integrate by parts $2N+2$ times
for $|\xi|>1$. This proves a remainder bound uniform over our
amplitudes, by \eqref{eq:ampbounds}. One may justify the Gaussian
Fourier interchange first with a factor $e^{-\varepsilon w^2}$ and
then let $\varepsilon\downarrow0$.

The actual amplitude here is a fixed cutoff times
\[
 e^{-\alpha d}E_\alpha(sY(w))Y'(w).
\]
At zero the cutoff is identically one. In the term indexed by $r$
in \eqref{eq:gaussian}, Taylor-expand $E_\alpha(sY(w))$ in $s$
through degree $N-r-1$. The retained terms are
\[
 \frac{s^k}{k!}E_\alpha^{(k)}(0)Y(w)^kY'(w),\qquad r+k<N.
\]
Real Taylor remainder estimates, also after the finitely many needed
$w$-derivatives, give a bound polynomial in $\alpha$ times
$e^{-c\alpha}s^{N-r}$. The damping remains valid on the real
Taylor segments $0\le\xi\le sY(w)$ because of
\eqref{eq:Wpositive}. Multiplying by
$s\lambda^{-r-1/2}$ gives exactly
$n^{-3/4-N/2}$ times a fixed power of $\alpha$.
The remainder in \eqref{eq:gaussianbound}, after the outside factor
$s$, has the same required form.

For the retained terms,
\[
 s\lambda^{-1/2}=(\alpha/\mu)^{1/4}n^{-3/4},\qquad
 s^k\lambda^{-r}
 =n^{-(r+k)/2}(\alpha/\mu)^{k/2}(\mu/\alpha)^{r/2}.
\]
These identities give precisely \eqref{eq:Qrecipe}. Since the
coefficients of $B$ are real, $W(-\theta)=\overline{W(\theta)}$.
The negative stationary point supplies the complex conjugate of the
positive one. The factor $2$ from adding them combines with
$1/(2\pi)$ and the Gaussian factor $\sqrt\pi$ to give
$1/\sqrt\pi$ in \eqref{eq:Jasym}.

\paragraph{The nonstationary complement.}
There are three ranges, each controlled by ordinary integrations by
parts with the phase derivative.
\begin{enumerate}
\item On a fixed compact set avoiding $0,1,-1$, $|\phi'|$ is bounded
below. After $L$ integrations the contribution, including the outside
factor $s$, is bounded by
$C_Ls\lambda^{-L}(1+\alpha)^{D_L}e^{-c\alpha}$.
\item On $0<|y|<1/2$, substitute $u=1/y$. The phase is again
$u+u^{-1}$ and the Jacobian is $u^{-2}$. The transformed amplitude and
every needed derivative have integrable tails as $|u|\to\infty$.
Boundary terms there vanish; at the other endpoint the partition
cutoff handles the boundary. The same bound follows.
\item On $|y|>2$, $\phi'$ is bounded away from zero and derivatives
of $1/\phi'$ decay integrably. For $k\ge1$,
\begin{equation}\label{eq:outerL1}
 \int\left|\frac{\mathrm d^k}{\mathrm dy^k}
       \{\chi(sy)e^{-\alpha W(sy)}\}\right|\dd y
 \le C_ks^{k-1}(1+\alpha)^{D_k}e^{-c\alpha}.
\end{equation}
This follows by changing variables back to $\theta$ and the chain
rule. Since $s\le s_0$, the bounds are uniform. Terms with no
derivative on the amplitude instead have an integrably decaying
derivative of $1/\phi'$. Thus $L\ge1$ integrations give the same
$C_Ls\lambda^{-L}$ bound. The outer smooth cutoff removes endpoint
terms at the support boundary.
\end{enumerate}
Take $L=N+1$. The power of $n$ in $s\lambda^{-L}$ is
$n^{-1-N/2}$, smaller than the target $n^{-3/4-N/2}$; powers of
$\alpha$ are absorbed in $(1+\alpha)^{M_N}$, since
$\alpha\ge\alpha_0$. This proves the lemma.
\end{proof}

\subsection{Summing actual mode errors}
Apply Lemma~\ref{lem:uniform} to $\alpha_m=2\pi^2m$ with
$\alpha_m\le\eta n$. The sum of the actual remainders is bounded by
\begin{equation}\label{eq:sumerror}
 C_Nn^{-3/4-N/2}
   \sum_{m\ge1}c_m(1+\alpha_m)^{M_N}e^{-c\alpha_m}
     =\OO_{B,N}(n^{-3/4-N/2}).
\end{equation}
For $\alpha_m>\eta n$, the integral itself satisfies
\[
 |J_{\alpha_m}(n)|\le C e^{-\alpha_md/2}.
\]
The entire large-mode tail is exponentially small in $n$. Each
retained coefficient series also has an exponentially small tail
there: its coefficient is an exponential times a fixed power of
$\alpha_m$. Therefore the partial sums of the retained terms may be
extended to all $m\ge1$. Finally, choose the smooth-circle remainder
order $L$ in \eqref{eq:reduction} large enough. This proves
\eqref{eq:general} with a true $n$-independent remainder constant.

\subsection{Reality and convergence of the coefficient recipe}
The symmetry $W(-\theta)=\overline{W(\theta)}$ implies that
$E_\alpha^{(k)}(0)$ is $(-\ii)^k$ times a real polynomial in
$\alpha$ of degree at most $k$. The derivatives involving $Y$ in
\eqref{eq:Qrecipe} are real. Thus $Q_j=(-\ii)^jP_j$. The factor
$\alpha^{(k-r)/2}=\alpha^{-j/2}\alpha^k$ shows that $P_j$ is
$\alpha^{-j/2}$ times a polynomial of degree at most $2j$.
Only finitely many cumulants enter at any fixed order.

The real part in \eqref{eq:Jasym} is
$P_j(\alpha)\cos(2u\sqrt{\alpha/\mu}+\pi/4+j\pi/2)$,
which gives \eqref{eq:Hrecipe} including its signs. Derivatives of
any fixed order in $u$ add only powers of $\sqrt{\alpha_m}$.
For complex $u$ in a compact set, the trigonometric factors have
absolute value at most a fixed multiple of
$\exp(2|\Ima u|\sqrt{\alpha_m/\mu})$. Exponential decay linear in
$m$ dominates this square-root growth and every polynomial.
Consequently all asserted convergences, analyticity and real-axis
boundedness follow. This completes the proof of
Theorem~\ref{thm:general}.

\section{Explicit terms for the logarithmic decoration}\label{sec:three}
For $B(z)=L(z)$, the solution $B(\rho)=1$ is
$\rho=1-e^{-1}<1$, strictly within its analytic disk. Its positive
coefficients occupy all sizes, so the theorem applies.
Using the logarithmic moment generating function
$\log L(\rho e^s)$ gives
\begin{equation}\label{eq:cumulants}
 \mu=v=M,\qquad
 \kappa_3=M^3+M,\qquad \kappa_4=M^4+6M^3+M.
\end{equation}
For instance $\mu=\rho/(1-\rho)=M$; the remaining identities follow
by successively applying $z\,\mathrm d/\mathrm dz$ to
$\log L(z)$ and evaluating at $z=\rho$. In particular,
\begin{equation}\label{eq:dqlog}
 d=\frac1{2M},\qquad q_1=\frac M6-\frac1{12M}.
\end{equation}
Define
\begin{equation}\label{eq:AKphi}
 A_m=\frac{2\pi^2m}{M},\qquad
 K_m=c_me^{-A_m/2}\frac{A_m^{1/4}}{\sqrt\pi},\qquad
 \Phi_m(u)=2u\sqrt{A_m}+\frac\pi4,
 \qquad U=\frac{2M^2-1}{12}.
\end{equation}
Then
\begin{align}
 H_0(u)&=-\sum_{m\ge1}K_m\cos\Phi_m(u),\label{eq:H0log}\\
 H_1(u)&=\sum_{m\ge1}K_m
     \frac{(8M^2-4)A_m^2+9}{48\sqrt{A_m}}\sin\Phi_m(u).
 \label{eq:H1log}
\end{align}
These expressions are absolutely convergent, and the tail decays
rapidly for this particular decoration.

\subsection{The third oscillatory coefficient with its sign}
For a general decoration, write
\[
 W(\theta)=d+\ii q_1\theta+r_2\theta^2+\OO(\theta^3).
\]
Reciprocal-series expansion of \eqref{eq:texpand} through the next
term yields
\begin{equation}\label{eq:r2}
 r_2=-\frac{v^3}{8\mu^4}
            +\frac{v\kappa_3}{6\mu^3}
            -\frac{\kappa_4}{24\mu^2}.
\end{equation}
The three terms with $r+k=2$ in \eqref{eq:Qrecipe} give
\begin{equation}\label{eq:Q2general}
 Q_2(\alpha)=-\frac{\alpha^3q_1^2}{2\mu}
          -\frac{\alpha^2r_2}{\mu}
          -\frac{15\alpha q_1}{16}
          +\frac{15\mu}{512\alpha}.
\end{equation}
To check the constants directly,
\[
 Y'(0)=1,\qquad (Y Y')''(0)=15/4,\qquad Y^{(5)}(0)=-15/16,
\]
while $E_\alpha'(0)=-\ii\alpha q_1$ and
$E_\alpha''(0)=-2\alpha r_2-\alpha^2q_1^2$.
Note that $Q_2=-P_2$, rather than $Q_2=P_2$.

For $B=L$, equations \eqref{eq:cumulants} and \eqref{eq:r2} simplify
to $r_2=-M^2/24-M/12$. To avoid confusing the scaled variable with
the universal recipe, denote $Q_2(MA)$ by $R_2(A)$. Then
\begin{equation}\label{eq:R2}
 R_2(A)=-\frac{A^3U^2}{2}
          +A^2\left(\frac{M^3}{24}+\frac{M^2}{12}\right)
          -\frac{15AU}{16}+\frac{15}{512A},
\end{equation}
and the third coefficient is
\begin{equation}\label{eq:H2log}
 H_2(u)=-\sum_{m\ge1}K_mR_2(A_m)\cos\Phi_m(u).
\end{equation}
Equations \eqref{eq:H0log}, \eqref{eq:H1log}, \eqref{eq:H2log},
and Theorem~\ref{thm:general} prove Theorem~\ref{thm:A}.

\subsection{An independent Laguerre route for finite diagnostics}
Put $w=1-z/\rho$. Direct series inversion gives
\begin{equation}\label{eq:localjet}
 \frac1{-\log L(z)}
   =\frac1{Mw}-\frac M6w+\frac{M^2}{24}w^2
                         -\frac{7M^3}{180}w^3+\OO(w^4).
\end{equation}
Thus, for a fixed dual parameter $\alpha$, the analytic amplitude
multiplying $\exp(-A/w)$, $A=\alpha/M$, begins as
\begin{equation}\label{eq:expjet}
 e^{-\alpha/[-\log L(z)]}
 =e^{-A/w}\left\{1+\frac{\alpha M}{6}w
       +M^2\left(\frac{\alpha^2}{72}-\frac\alpha{24}\right)w^2
       +\OO_\alpha(w^3)\right\}.
\end{equation}
The associated-Laguerre generating function implies the exact identity
\begin{equation}\label{eq:laguerre}
 [x^n](1-x)^j e^{-A/(1-x)}=e^{-A}L_n^{(-j-1)}(A).
\end{equation}
It follows simply by multiplying
$\sum_{n\ge0}L_n^{(\nu)}(A)x^n=(1-x)^{-\nu-1}
\exp(-Ax/(1-x))$ by $e^{-A}$ and taking $\nu=-j-1$.
Finite jets in \eqref{eq:expjet} therefore give independently
computable Laguerre approximants.

The fixed-positive-argument Laguerre expansions in
DLMF 18.15(iv)~\cite{laguerre} agree with the first terms above.
Their stated compact-argument uniformity does not cover a dual
parameter tending to infinity. Accordingly, this is a cross-check and
numerical route; it is not used in place of the uniform summation
argument of Section~\ref{sec:uniform}.

\section{A rational tail bound and genuine sign changes}\label{sec:sign}
For the logarithmic decoration, put
\[
 r=e^{-\pi^2/M},\qquad
 K_1=e^{-\pi^2/M}\frac{(2\pi^2/M)^{1/4}}{\sqrt\pi}>0.
\]
Numerically $r=0.003202319565807\ldots$ and
$K_1=0.0033262016024504409820\ldots$. We next prove the bound needed
for sign changes using only rational inequalities, without relying on
those decimals.

\begin{lemma}[First-harmonic dominance]\label{lem:tail}
For $K_m$ defined in \eqref{eq:AKphi},
\begin{equation}\label{eq:tail}
 \frac{\sum_{m\ge2}K_m}{K_1}
 <\frac{359101}{26730899}<\frac1{50}.
\end{equation}
\end{lemma}
\begin{proof}
An elementary upper bound for $e$ is
\[
 e<\sum_{k=0}^6\frac1{k!}
           +\frac{1/7!}{1-1/8}
     =\frac{31967}{11760}<\frac{87}{32}.
\]
The geometric bound is valid since successive terms after $1/7!$
have ratios at most $1/8$. Hence $M<55/32$. Archimedes' classical
bound $\pi>223/71>157/50$ gives
\[
 \frac{\pi^2}{M}>\frac{(157/50)^2}{55/32}
     =\frac{197192}{34375}>\frac{573}{100}.
\]
The finite positive rational sum
$\sum_{k=0}^{20}(573/100)^k/k!$ is greater than $300$.
Therefore $e^{\pi^2/M}>300$, and $r<1/300$.

By Lemma~\ref{lem:cm},
\[
 \frac{K_m}{K_1}=c_m m^{1/4}r^{m-1}
        \le(1+\log m)m^{1/4}r^{m-1}
        \le m^2r^{m-1}\quad(m\ge2).
\]
For the last inequality, $1+\log m\le m$ and $m^{1/4}\le m$
suffice. Summing $\sum_{m\ge1}m^2r^{m-1}=(1+r)/(1-r)^3$ and
using $r<1/300$ yields
\[
 \sum_{m\ge2}\frac{K_m}{K_1}
 <\frac{1+1/300}{(1-1/300)^3}-1
 =\frac{359101}{26730899}<\frac1{50}.
\]
All the final comparisons are rational and are replayed by the exact
verification program.
\end{proof}

\begin{theorem}[Signed subsequences]\label{thm:sign}
For A330499,
\begin{align}
 \liminf_{n\to\infty}n^{3/4}
       \left(\frac{a_n\rho^n}{n!}-C\right)&\le-\frac{49K_1}{50},\label{eq:liminf}\\
 \limsup_{n\to\infty}n^{3/4}
       \left(\frac{a_n\rho^n}{n!}-C\right)&\ge\frac{49K_1}{50}.
 \label{eq:limsup}
\end{align}
In particular the normalized error is not $o(n^{-3/4})$ and admits
no pure integer-power refinement beginning with $C+c_1/n$.
\end{theorem}
\begin{proof}
Set $A_1=2\pi^2/M$. For all sufficiently large integers $j$, choose
$n_j$ to be a nearest integer to
\[
 \left(\frac{2\pi j-\pi/4}{2\sqrt{A_1}}\right)^2.
\]
Changing this square by at most $1/2$ changes its square root by
$\OO(1/j)$. Hence $\cos\Phi_1(\sqrt{n_j})\to1$. Equation
\eqref{eq:H0log} and Lemma~\ref{lem:tail} imply
$\limsup_j H_0(\sqrt{n_j})\le-49K_1/50$.
Choosing instead a nearest integer to
$(((2j+1)\pi-\pi/4)/(2\sqrt{A_1}))^2$ makes the first cosine tend
to $-1$, giving the opposite bound. These are increasing unbounded
subsequences after discarding finitely many initial values.
Finally, Theorem~\ref{thm:A} with $N=1$ gives
\[
 n^{3/4}\left(\frac{a_n\rho^n}{n!}-C\right)
          =H_0(\sqrt n)+\OO(n^{-1/2}).
\]
This proves both comparisons. A refinement $C+c_1/n+\OO(n^{-2})$
would make this expression tend to zero, a contradiction.
\end{proof}
The strict tail inequality actually permits strict improvements of
\eqref{eq:liminf}--\eqref{eq:limsup}. No exact value of either extreme
is claimed. Their determination would require controlling the joint
phase behavior, not just the dominant harmonic.

\section{Actual root of unity cusps and the Delta obstruction}\label{sec:cusps}
The failure of a usual analytic-continuation hypothesis should not be
inferred just from the appearance of an infinite product. The
following argument proves genuine singularities directly.

\begin{proposition}[Odd-order cusps]\label{prop:cusps}
Let $\zeta$ be a primitive root of unity of odd order $h$. With the
analytic logarithmic branch of $F_0$ in the unit disk,
\begin{equation}\label{eq:cusp}
 \lim_{t\downarrow0}tF_0(\zeta e^{-t})=\frac{\pi^2}{12h^2}.
\end{equation}
In particular each such root is a nonremovable singularity.
\end{proposition}
\begin{proof}
Absolute convergence for $|q|<1$ permits the Lambert representation
\[
 \log P(q)=-\sum_{r\ge1}\frac{q^r}{r(1-q^r)}.
\]
At $q=\zeta e^{-t}$ the terms $r=h\ell$ contribute, after
multiplication by $t$,
\[
 -\sum_{\ell\ge1}\frac{t}{h\ell(e^{h\ell t}-1)}
       \longrightarrow -\frac{\pi^2}{6h^2}.
\]
Dominated convergence is justified by
$t/(e^{h\ell t}-1)\le1/(h\ell)$.
For $r$ not divisible by $h$, the finitely many nontrivial root
arguments have radial segments bounded away from $1$. Thus
$|1-\zeta^re^{-rt}|\ge c_h>0$, uniformly in $r,t$. The absolute
contribution of these terms is at most
\[
 c_h^{-1}\sum_{r\ge1}\frac{e^{-rt}}r=\OO_h(\log(1/t)),
\]
which disappears after multiplication by $t$.
Since $h$ is odd, $\zeta^2$ still has order $h$. Applying the same
limit to $q^2=\zeta^2e^{-2t}$ yields
$t\log P(q^2)\to-\pi^2/(12h^2)$.
Subtracting $t\log P(q)$ proves \eqref{eq:cusp}.
The order $h$ is fixed before taking the limit; no uniform-in-$h$
estimate has been used.
\end{proof}

\begin{corollary}[No standard Delta-domain at the positive point]\label{cor:delta}
For any decoration in Theorem~\ref{thm:general}, $F_B$ has no
holomorphic continuation to a standard Delta-neighborhood of $\rho$
that excludes only an outward wedge of half-angle less than $\pi/2$.
\end{corollary}
\begin{proof}
Since $B'(\rho)=\mu/\rho>0$, let $g$ be its local holomorphic
inverse near $B(\rho)=1$. For odd $h\to\infty$, put
$\zeta_h=e^{2\pi\ii/h}$ and $z_h=g(\zeta_h)$. Taylor expansion gives
\[
 z_h=\rho+\frac\rho\mu(\zeta_h-1)+\OO(|\zeta_h-1|^2),
 \quad z_h\to\rho,\quad \arg(z_h-\rho)\to\pi/2.
\]
The conjugate sequence has limiting argument $-\pi/2$.
The local inverse image of the disk $|q|<1$ is connected near
$\rho$ and meets the original coefficient disk. Along the radial
curves $g(\zeta_he^{-t})$, Proposition~\ref{prop:cusps} gives a
divergence at $z_h$. This is the same analytic branch as $F_B$ by
continuation from their common interior region.

A standard Delta-neighborhood consists locally of points sufficiently
close to $\rho$, except for an outward wedge
$|\arg(z-\rho)|\le\phi$ with some $\phi<\pi/2$.
For all sufficiently large $h$, $z_h$ lies strictly inside the
remaining neighborhood. The nearby radial curves also lie there and
connect to the interior branch through the local inverse image of
$|q|<1$. A holomorphic continuation to that whole neighborhood would
therefore have to be regular at $z_h$, contradicting the divergence.
\end{proof}

This establishes the obstruction needed to select the proof method.
It does not classify all root-of-unity cusps or all possible analytic
continuation domains. The boundary proof is independently valid and
does not use the obstruction as a substitute for an error estimate.

\section{Smooth inverses and pointwise integer envelopes}\label{sec:inverse}
The coefficient theorem yields arbitrarily accurate fixed-order
\emph{smooth charts}. The discrete threshold remains a staircase, so
its statement must retain rounding explicitly. Throughout this
section the sequence is A330499 and the constants are
\eqref{eq:logconstants}.

\subsection{Every fixed order with two ceilings}
Define, for $Y>0$,
\begin{equation}\label{eq:threshold}
 \mathcal N(Y)=\min\{n\ge1:a_n\ge Y\}.
\end{equation}
Proposition~\ref{prop:counts} proves strict monotonicity, and
Theorem~\ref{thm:A} proves unboundedness, so this is well-defined.
For a fixed integer $J\ge0$, put
\begin{align}
 S_J(x)&=C+\sum_{j=0}^Jx^{-3/4-j/2}H_j(\sqrt x),\label{eq:SJ}\\
 \mathcal L_J(x)&=\log\Gamma(x+1)-x\log\rho+\log S_J(x),\label{eq:LJ}\\
 p_J&=\frac34+\frac{J+1}{2}.\label{eq:pJ}
\end{align}
These functions are defined for all sufficiently large real $x$,
since $S_J(x)=C+\OO(x^{-3/4})>0$.

\begin{theorem}[Pointwise threshold envelopes]\label{thm:inverse}
For each fixed $J\ge0$, there are $D_J>0$ and a starting point
$x_J$ such that the charts
$\mathcal L_J(x)\pm D_Jx^{-p_J}$ are strictly increasing on
$x\ge x_J$, and their large roots defined by
\begin{align}
 \mathcal L_J(x_+)+D_Jx_+^{-p_J}&=\log Y,\label{eq:xplus}\\
 \mathcal L_J(x_-)-D_Jx_-^{-p_J}&=\log Y\label{eq:xminus}
\end{align}
satisfy, for every sufficiently large $Y$,
\begin{equation}\label{eq:ceil}
 \lceil x_+\rceil\le\mathcal N(Y)\le\lceil x_-\rceil.
\end{equation}
Moreover $x_+\le x_-$ and
\begin{equation}\label{eq:gap}
 x_--x_+=\OO_J\left(\frac{x_+^{-p_J}}{\log x_+}\right).
\end{equation}
The constants and starting point are existence statements, not
numerically certified values.
\end{theorem}
\begin{proof}
Theorem~\ref{thm:A} with $N=J+1$, followed by the logarithm on a
fixed neighborhood of the positive constant $C$, gives
\begin{equation}\label{eq:logbackerror}
 |\log a_n-\mathcal L_J(n)|\le D_Jn^{-p_J}
\end{equation}
for all sufficiently large integers $n$, after choosing $D_J$.
Boundedness of the derivatives of $H_j$ gives
$S_J'(x)=\OO_J(x^{-5/4})$. The standard digamma expansion implies
\[
 \mathcal L_J'(x)
 =\psi(x+1)-\log\rho+\frac{S_J'(x)}{S_J(x)}
 =\log(x/\rho)+\OO(1/x)+\OO_J(x^{-5/4}).
\]
Both perturbed charts therefore have positive derivative for large
$x$ and tend to infinity. Their roots exist uniquely there, and the
upper chart reaches a given height first, so $x_+\le x_-$.

If an integer $n<x_+$ is sufficiently large, its upper chart is
less than $\log Y$, and \eqref{eq:logbackerror} forces $a_n<Y$.
The finitely many smaller integers also satisfy this once $Y$ is
large enough. Hence $\mathcal N(Y)\ge\lceil x_+\rceil$.
For $n=\lceil x_-\rceil$, the lower chart is at least $\log Y$,
so $a_n\ge Y$ and the upper bound follows. This reasoning remains
valid if either root is an integer, or if $Y$ is arbitrarily close to
a sequence value.

Subtracting the root equations and using the lower derivative bound
$c\log x$ for $\mathcal L_J$ gives \eqref{eq:gap} by the mean value
theorem. The roots are asymptotically equivalent, as follows already
from their common leading factorial scale.
\end{proof}

For large $Y$ the two ceilings are equal or adjacent, since the real
root gap tends to zero. They cannot uniformly be collapsed to one
ceiling of the central chart. Arbitrarily small changes in $Y$ near
a jump can change $\mathcal N(Y)$, however accurate the smooth chart
may be.

\subsection{The first oscillatory displacement of the exact core}
Let $x_*=x_*(Y)$ be the unique sufficiently large root of the
\emph{exact factorial core}
\begin{equation}\label{eq:core}
 \mathcal F_*(x)=\log\Gamma(x+1)-x\log\rho+\log C=\log Y.
\end{equation}
Its derivative is
$\mathcal F_*'(x)=\psi(x+1)-\log\rho\sim\log(x/\rho)$.

\begin{proposition}[First oscillatory inverse shift]\label{prop:shift}
Set
\begin{equation}\label{eq:shift}
 x_{\mathrm{osc}}=x_*-
 \frac{x_*^{-3/4}H_0(\sqrt{x_*})}
      {C\{\psi(x_*+1)-\log\rho\}}.
\end{equation}
The central smooth root of $\mathcal L_0(x)=\log Y$ differs from
$x_{\mathrm{osc}}$ by
$\OO(x_*^{-3/2}/\log x_*)$. In particular there is a fixed $K>0$
such that the actual integer threshold satisfies
\begin{equation}\label{eq:firstceil}
 \left\lceil x_{\mathrm{osc}}-
            \frac{Kx_*^{-5/4}}{\log x_*}\right\rceil
 \le\mathcal N(Y)\le
 \left\lceil x_{\mathrm{osc}}+
            \frac{Kx_*^{-5/4}}{\log x_*}\right\rceil
\end{equation}
for all sufficiently large $Y$.
\end{proposition}
\begin{proof}
Write $E_0(x)=\log(S_0(x)/C)$. Then
\[
 E_0(x)=C^{-1}x^{-3/4}H_0(\sqrt x)+\OO(x^{-3/2}),
 \qquad E_0'(x)=\OO(x^{-5/4}).
\]
The central root of $\mathcal F_*+E_0$ differs from $x_*$ by
$\OO(x_*^{-3/4}/\log x_*)$. Taylor expansion over that interval
shows that its first displacement is
$-E_0(x_*)/\mathcal F_*'(x_*)$, with an error bounded by
$\OO(x_*^{-2}/(\log x_*)^2)$ from varying $E_0$ and smaller
core-curvature terms. Replacing $E_0(x_*)$ by its first displayed
term leaves the stated $\OO(x_*^{-3/2}/\log x_*)$ error.
The two roots of Theorem~\ref{thm:inverse} for $J=0$ are within
$\OO(x_*^{-5/4}/\log x_*)$ of the central root. Enlarge a single
constant $K$ to contain all these errors and take ceilings.
\end{proof}

\subsection{Lambert initialization and the accuracy it does not supply}
Write $y=\log Y$ and use the positive real branch of Lambert $W$.
The solution of $x_0\log(x_0/(e\rho))=y$ is
\begin{equation}\label{eq:lambert}
 x_0=\frac{y}{W(y/(e\rho))}.
\end{equation}
Stirling's expansion gives
\[
 \mathcal F_*(x)=x\log\frac{x}{e\rho}
        +\frac12\log(2\pi x)+\log C+\frac1{12x}+\OO(x^{-3}).
\]
Since the omitted half-logarithmic term is $\OO(\log x_0)$ and the
core derivative is of order $\log x_0$, the exact core root first
differs from $x_0$ by $\OO(1)$. One Taylor correction yields
\begin{equation}\label{eq:lambertcorrect}
 x_*=x_0-
       \frac{\frac12\log(2\pi x_0)+\log C}{\log(x_0/\rho)}
        +\OO\left(\frac1{x_0\log x_0}\right).
\end{equation}
Indeed the remaining Stirling term, derivative variation and quadratic
core term are all $\OO(1/x_0)$ in the equation, hence
$\OO(1/(x_0\log x_0))$ in the root.

This Lambert expression is a useful seed for solving
\eqref{eq:core}. Its displayed error is smaller than the leading
oscillatory displacement, but larger than the next half-power inverse
corrections. Therefore replacing $x_*$ in an arbitrarily high-order
claim by this finite seed without propagating its error would be
incorrect. The root in \eqref{eq:shift} is the exact factorial-core
root, not just \eqref{eq:lambertcorrect}.

\subsection{Generating higher smooth inverse terms}
For any fixed $J$, put $E_J(x)=\log(S_J(x)/C)$, so that
$\mathcal L_J=\mathcal F_*+E_J$. Let $X(y)$ be the large local
inverse of $\mathcal F_*$ and write $D_y=(1/\mathcal F_*'(x))D_x$
when acting on a function of $x=X(y)$. Ordinary Lagrange reversion
of the perturbation parameter $\varepsilon$ in
$\mathcal F_*(x)+\varepsilon E_J(x)=y$ gives
\begin{equation}\label{eq:reversion}
 x=X(y)+\sum_{k\ge1}\frac{(-\varepsilon)^k}{k!}
     D_y^{k-1}\left\{X'(y)E_J(X(y))^k\right\}.
\end{equation}
This is a local analytic identity in $\varepsilon$ near zero and a
constructive finite-order reversion rule. For sufficiently large
$y$, one can set $\varepsilon=1$: in a fixed complex neighborhood
of the real root the perturbation and its first derivative tend to
zero, whereas the core derivative has size $\log X(y)$, so the
analytic implicit solution exists on a parameter disk containing
$|\varepsilon|\le1$. Equivalently, finite Taylor reversion with a
remainder follows by repeated substitution and the same derivative
bounds.

Only finitely many derivatives and products are needed for any
requested fixed precision. Increasing $J$ reduces the coefficient
backward error; reversion then converts it into a smooth root error
of order $x^{-p_J}/\log x$. Finally, Theorem~\ref{thm:inverse}
converts that information to a pointwise integer statement. These are
classical inverse-chart and backward-error operations; the
model-specific input is the oscillatory coefficient law. They do not
create an unambiguous real inverse of the discrete staircase.

\section{Computational checks and reproducibility}\label{sec:computations}
The companion archive separates mathematical proof, exact finite
checks and floating-point diagnostics. Agreement at many indices is
useful for detecting algebraic errors, but is not a replacement for
Lemma~\ref{lem:uniform}, nor an interval certificate for its implicit
constant.

\subsection{Exact mandatory checks}
The mandatory Python verification kernel uses only the standard
library. It computes the positive Stirling transform
\eqref{eq:stirling}, matches all 21 terms listed in the inspected
OEIS record, checks the direct alternating-divisor identity through
$100$, checks the exact recurrence and strict growth through $600$,
and replays the rational inequalities in Lemma~\ref{lem:tail}.
Independent finite coefficient computations test the formal
composition directly rather than reusing the same weighted transform.

Checks use explicit exception-raising guards, so Python's
optimization flag cannot silently disable them. Negative controls
alter fixtures and other mathematical inputs and must fail in both
ordinary Python and Python \texttt{-O}. A fixed checksum of the
exact integer output through $600$ is a reproducibility fingerprint;
it is not evidence in place of the underlying identities.

\subsection{Optional floating diagnostics}
The optional high-precision checks use \texttt{mpmath}; the FFT and
two-size experiment additionally use the numerical dependencies
listed in the archive. They are excluded from the standard-library
mathematical gate. The portable diagnostic scripts and the historical numerical
outputs are retained with provenance and hashes, and the high-precision
replay compares 50- and 80-digit computations. Finite tolerances are
regression limits on those samples, not uniform asymptotic bounds.

For numerical comparison define
\[
 T_n=n^{3/4}\left(\frac{a_n\rho^n}{n!}-C\right),\qquad
 T_n^{(J)}=\sum_{j=0}^Jn^{-j/2}H_j(\sqrt n).
\]
The first two error columns in Table~\ref{tab:coeffdiag} illustrate that
$T_n$ can remain visibly displaced from $H_0(\sqrt n)$ at moderate
indices. Retaining $H_1$ generally helps, but asymptotic truncation
does not promise monotone improvement at every individual $n$.
\begin{table}[htbp]
\centering
\small
\begin{tabular}{r r r r}
\toprule
$n$ & $T_n$ & $T_n^{(0)}$ & $T_n^{(1)}$\\
\midrule
100   & $-0.003509157374$ & $-0.002844061648$ & $-0.005567577555$\\
200   & $ 0.004071822395$ & $ 0.002461800049$ & $ 0.004970742765$\\
400   & $-0.001726451607$ & $ 0.000985440349$ & $-0.001563360516$\\
800   & $ 0.000498855063$ & $ 0.002121973187$ & $ 0.000685065706$\\
1000  & $ 0.001520405923$ & $-0.000167680657$ & $ 0.001516496146$\\
2000  & $ 0.003100045820$ & $ 0.002332031756$ & $ 0.003185266147$\\
\bottomrule
\end{tabular}
\caption{Finite numerical diagnostics, rounded for display.
Exact integer coefficients underlie $T_n$; harmonic sums are numerical.
No entry is an interval enclosure.}\label{tab:coeffdiag}
\end{table}

For inverse diagnostics one may take $Y=a_n$, whose exact threshold
is $n$, and compare smooth roots with that known integer.
Table~\ref{tab:inversediag} uses the exact factorial core, the first
shift \eqref{eq:shift}, and the root of the three-term chart
$\mathcal L_2$. The small residuals illustrate the formulas only;
they do not license one-ceiling rounding near every threshold.
\begin{table}[htbp]
\centering
\small
\begin{tabular}{r r r r}
\toprule
$n$ & $x_*-n$ & $x_{\mathrm{osc}}-n$ & $\mathcal L_2^{-1}(\log a_n)-n$\\
\midrule
100  & $-4.57425935\,10^{-5}$ & $-8.67428957\,10^{-6}$ & $ 1.05433425\,10^{-6}$\\
200  & $ 2.77698434\,10^{-5}$ & $ 1.09789162\,10^{-5}$ & $-1.46614723\,10^{-7}$\\
400  & $-6.25087523\,10^{-6}$ & $-9.81874913\,10^{-6}$ & $ 6.74257110\,10^{-8}$\\
800  & $ 9.69829090\,10^{-7}$ & $-3.15553380\,10^{-6}$ & $ 4.07413144\,10^{-9}$\\
1000 & $ 2.42462740\,10^{-6}$ & $ 2.69203273\,10^{-6}$ & $-7.70215372\,10^{-9}$\\
2000 & $ 2.68683539\,10^{-6}$ & $ 6.65621847\,10^{-7}$ & $-1.17990530\,10^{-9}$\\
\bottomrule
\end{tabular}
\caption{Floating-point smooth inverse errors at known jump values.
A first correction need not improve every finite case.}\label{tab:inversediag}
\end{table}

Two independent diagnostics are especially useful for detecting
model-specific mistakes. The two-size decoration has the finite
coefficient formula
\begin{equation}\label{eq:twosize}
 [z^n]F_{(z+z^2)/2}(z)
   =\sum_{k=\lceil n/2\rceil}^{n}
         c_k2^{-k}\binom{k}{n-k}.
\end{equation}
Here $\rho=1$, $\mu=3/2$, $v=1/4$, $\kappa_3=0$ and
$q_1=-1/216$. Comparisons through $n=20000$ test the universal
first correction against a coefficient formula unrelated to
Stirling numbers. Separately, a modular-transform FFT with
$262144$ samples and $80$ eta modes extracts coefficients through
selected indices up to $64000$, with overlap checks against exact
integers through $2000$. These FFT outputs are ordinary floating-point
experiments, not validated quadrature or interval calculations.

\subsection{Rebuilding and the integrity boundary}
The archive contains the editable \LaTeX{} source, the PDF, the exact
kernel, optional experiments, machine-readable check results, a
manifest and rebuilding instructions. The build uses a fresh working
directory, rejects relevant \LaTeX{} warnings and layout overflows,
and writes a deterministically ordered ZIP with normalized metadata.
The release is checked by clean extraction and rebuilding; changing a
hashed payload must be rejected by its integrity checker.

These guarantees concern the specified package and toolchain.
Different \TeX{} or font versions may change PDF bytes while preserving
mathematical content. Full analytic proofs are in this manuscript;
no third-party papers or private research notes are needed for a
rebuild. See the archive's reproducibility guide for the exact commands,
optional dependencies and the distinction between the exact and
numerical test tiers.

\section{Source comparison and attribution}\label{sec:sources}
The public sources below were reviewed for the mathematical model and
proof inputs. The comparison is deliberately bounded.

\paragraph{OEIS model.}
The official A330499 record~\cite{oeis}, revision 11 dated
17 December 2019, defines \eqref{eq:model} and lists the leading
numerical constant. Its inspected formula field does not give the
closed form \eqref{eq:logconstants}, a proof, or higher corrections.
The adjacent A330498 has an additional harmonic $1/k$ weight in its
kernel. It does not satisfy the eta identity used here and is excluded
from every theorem in this report.

\paragraph{Eta transformation.}
The modular identity is classical, as recorded in DLMF
23.18~\cite{eta}. Ardehali--Rosengren's 2026 paper~\cite{ardehali}
provides an explicit product formulation in equation (15). Its
small-parameter asymptotic discussion fixes a nontangential angular
range in the right half-plane. The present boundary
$t(\theta)=-\ii\mu\theta+v\theta^2/2+\cdots$ is tangential.
That source's displayed radial-type expansions therefore do not,
by themselves, prove the decorated coefficient theorem. Our use of
the exact identity is followed by a separate uniform boundary argument.

\paragraph{Classical coefficient and inverse machinery.}
The supernecklace construction is recorded in OEIS A003713 and the
Flajolet--Sedgewick book companion, exercise IV.28~\cite{supernecklace,ac}.
DLMF 18.15(iv)~\cite{laguerre} supplies classical Laguerre
asymptotics, with its stated compact-positive-argument uniformity.
The stationary-phase Gaussian expansion, Lambert factorial
initialization, local Lagrange reversion and backward-error conversion
are standard methods. The specific contribution here is their
uniform, summable application to a nonmonomial analytic decoration
and the ensuing explicit oscillatory laws and threshold envelopes.

\paragraph{Bounded overlap screen.}
Exact-sequence and broader model searches in the available report
collection and mathematical repository found no inspected report
with this A330499 decorated-logarithmic-product theorem. Potentially
adjacent inspected models involved critical polylogarithmic Lambert
charts, divisor-weighted exponentiated products, planar-map
coefficients on a different logarithmic scale, or factorially forced
recurrences. They did not supply the present boundary transfer.
This records the scope of the comparison, not proof of complete
private-collection coverage or worldwide novelty. No external
repository or OEIS record is changed by this report.

\section{Further questions and present limits}\label{sec:further}
The theorem gives every fixed order for a fixed decoration. Several
natural extensions require additional work.
\begin{enumerate}
\item \textbf{Effective error constants and cutoffs.}
The proof tracks polynomial-times-exponential bounds but does not
optimize explicit cutoff radii, amplitude norms or stationary-phase
constants. An interval implementation of these choices could turn
Theorem~\ref{thm:inverse} into a certified inverse algorithm. The
current finite diagnostics do not supply those missing constants.
\item \textbf{Other support spans.}
If the difference span exceeds one, multiple unit-modulus boundary
contacts occur. Their phases and cusp types must be combined. Merely
replacing the span-one condition by gcd of the occupied sizes is
invalid, as $(z+z^3)/2$ demonstrates.
\item \textbf{Vanishing variance.}
As $v/\mu^2$ decreases, the number of appreciable harmonics grows.
A uniform crossover theorem toward monomial decorations would need
new two-parameter estimates. It is not a consequence of keeping
$B$ fixed in Theorem~\ref{thm:general}.
\item \textbf{Joint phases and limiting error laws.}
The frequencies are proportional to $\sqrt m$. Their arithmetic
relations may permit sharper extreme-value or distribution results
for $H_0(\sqrt n)$. The signed subsequences here need only the first
harmonic and establish neither an empirical limit law nor the exact
limsup and liminf.
\item \textbf{Growing orders and exponentially improved expansions.}
The coefficient recipe is finite at every fixed order. It does not
control coefficient growth as the requested order tends to infinity,
optimal truncation, or a complete exponentially improved transseries.
\item \textbf{Weighted kernels.}
The harmonic weighting in A330498 leads to a different transformation
problem. Its neighboring sequence number is not evidence that the
same theorem applies. Other products or weights should be treated
from their own exact kernels.
\end{enumerate}

\paragraph{Conclusion.}
For A330499, the exact leading constant is only the first part of the
answer. Analytic decoration damps the eta modes by a variance ratio
but leaves a real $n^{-3/4}$ oscillation, followed by explicitly
computable half-power orders. A uniform boundary proof makes that
infinite oscillatory expansion rigorous and gives inverse envelopes
that remain valid at the jumps of the integer threshold.

\begin{thebibliography}{9}
\bibitem{oeis}
The OEIS Foundation Inc., \emph{A330499}, contributed by V. Kote\v{s}ovec,
16 December 2019; inspected revision 11, 17 December 2019.
\url{https://oeis.org/A330499}. Accessed 3 October 2026.
\bibitem{eta}
NIST Digital Library of Mathematical Functions, \emph{Section 23.18,
Modular Transformations}, especially equations 23.18.5--23.18.7.
\url{https://dlmf.nist.gov/23.18}. Accessed 3 October 2026.
\bibitem{ardehali}
A. A. Ardehali and H. Rosengren, \emph{A new product formula for
$(z;q)_\infty$, with applications to asymptotics}, arXiv:2602.11329v1,
2026; equation (15) and the angular scope of Section 4.
\url{https://arxiv.org/html/2602.11329v1}.
\bibitem{laguerre}
NIST Digital Library of Mathematical Functions, \emph{Section 18.15(iv),
Laguerre polynomials}, especially equations 18.15.14--18.15.16.
\url{https://dlmf.nist.gov/18.15}. Accessed 3 October 2026.
\bibitem{supernecklace}
The OEIS Foundation Inc., \emph{A003713, Supernecklaces}.
\url{https://oeis.org/A003713}. Accessed 3 October 2026.
\bibitem{ac}
P. Flajolet and R. Sedgewick, \emph{Analytic Combinatorics}, Cambridge
University Press, 2009; book companion, Chapter IV, exercise IV.28.
\url{https://ac.cs.princeton.edu/40complex/}.
\end{thebibliography}
\end{document}
